\subsection{Sucesiones espectrales de asociatividad de los hipertor globales}
\label{subsection:III.6.8}
\label{III.6.8}

\begin{env}[6.8.1]
\label{III.6.8.1}
Mantengamos las hipótesis y notaciones de \sref{III.6.7.11}; en
particular, los complejos $\sh P_\bullet^{(i)}$ están acotados
inferiormente. Sea $(I_j)_{j\in J}$ una partición del conjunto de
índices $I$. Buscamos una relación de «asociatividad» entre
\[
  \mathfrak{Tor}_n^S
  \bigl(
    (f_i)_{i\in I};
    (\sh P_\bullet^{(i)})_{i\in I}
  \bigr)
\]
y los hipertor «parciales»
\[
  \mathcal T_{\bullet j}
  =
  \mathfrak{Tor}_\bullet^S
  \bigl(
    (f_i)_{i\in I_j};
    (\sh P_\bullet^{(i)})_{i\in I_j}
  \bigr).
\]
\oldpage[III]{165}

Para simplificar la notación, nos limitaremos al caso en que $I$
es el intervalo $1\leq i\leq m$ y la partición consta de los dos
intervalos
\[
  \{1,2,\ldots,r\}
  \qquad\hbox{y}\qquad
  \{r+1,\ldots,m\}.
\]
\end{env}

\begin{proposition}[6.8.2]
\label{III.6.8.2}
Existe un funtor espectral canónico birregular, denominado
\emph{funtor espectral de asociatividad}, que se denota por
\[
  {}^{(e)}\sh E^S
  \bigl(
    f_1,\ldots,f_m;
    \sh P_\bullet^{(1)},\ldots,\sh P_\bullet^{(m)}
  \bigr),
\]
o simplemente por ${}^{(e)}\sh E$ con los mismos argumentos, que
converge a
\[
  \mathfrak{Tor}_\bullet^S
  \bigl(
    f_1,\ldots,f_m;
    \sh P_\bullet^{(1)},\ldots,\sh P_\bullet^{(m)}
  \bigr)
\]
y cuyo término $E_2$ es
\[
\begin{split}
  {}^{(e)}\sh E_{pq}^2
  =
  \bigoplus_{q_1+q_2=q}
  \shTor_p^S
  \Bigl(
    &
    \mathfrak{Tor}_{q_1}^S
    \bigl(
      f_1,\ldots,f_r;
      \sh P_\bullet^{(1)},\ldots,\sh P_\bullet^{(r)}
    \bigr),
    \\
    &
    \mathfrak{Tor}_{q_2}^S
    \bigl(
      f_{r+1},\ldots,f_m;
      \sh P_\bullet^{(r+1)},\ldots,\sh P_\bullet^{(m)}
    \bigr)
  \Bigr).
\end{split}
\]
\end{proposition}

\begin{proof}
Aquí se identifica canónicamente $Y$ con
$Z^{(1)}\times_SZ^{(2)}$, donde
\[
  Z^{(1)}
  =
  Y^{(1)}\times_S\cdots\times_SY^{(r)}
  \quad\hbox{y}\quad
  Z^{(2)}
  =
  Y^{(r+1)}\times_S\cdots\times_SY^{(m)}.
\]
Basta tratar el caso en que $S$ y todos los $Y^{(i)}$ son afines.
El paso de este caso al caso general utiliza los métodos de
\sref{III.6.7.1} y \sref{III.6.7.2}; dejamos los detalles al
lector. Basta, pues, demostrar el siguiente enunciado afín.
\end{proof}

\begin{corollary}[6.8.3]
\label{III.6.8.3}
Sea $A$ un anillo, sea $S=\Spec(A)$ y sean $X^{(i)}$,
$1\leq i\leq m$, $S$-esquemas cuasicompactos. Para cada $i$, sea
$\sh P_\bullet^{(i)}$ un complejo acotado inferiormente de
$\sh O_{X^{(i)}}$-módulos cuasicoherentes. Existe un funtor
espectral canónico birregular que converge a
\[
  \mathbf{Tor}_\bullet^S
  \bigl(
    X^{(1)},\ldots,X^{(m)};
    \sh P_\bullet^{(1)},\ldots,\sh P_\bullet^{(m)}
  \bigr)
\]
y cuyo término $E_2$ es
\[
\begin{split}
  {}^{(e)}E_{pq}^2
  =
  \bigoplus_{q_1+q_2=q}
  \Tor_p^A
  \Bigl(
    &
    \mathbf{Tor}_{q_1}^S
    \bigl(
      X^{(1)},\ldots,X^{(r)};
      \sh P_\bullet^{(1)},\ldots,\sh P_\bullet^{(r)}
    \bigr),
    \\
    &
    \mathbf{Tor}_{q_2}^S
    \bigl(
      X^{(r+1)},\ldots,X^{(m)};
      \sh P_\bullet^{(r+1)},\ldots,\sh P_\bullet^{(m)}
    \bigr)
  \Bigr).
\end{split}
\]
\end{corollary}

\begin{proof}
Siguiendo la definición de \sref{III.6.6.2}, elijamos para cada
$i$ un recubrimiento abierto afín finito $\mathfrak U^{(i)}$ de
$X^{(i)}$, formemos el bicomplejo
\[
  L_{\bullet\bullet}^{(i)}
  =
  C^\bullet
  \bigl(
    \mathfrak U^{(i)},\sh P_\bullet^{(i)}
  \bigr),
\]
y elijamos una resolución proyectiva de Cartan--Eilenberg
$M_{\bullet\bullet\bullet}^{(i)}$ de este bicomplejo en el sentido
de \sref[0]{0.11.7.1}. El hipertor considerado es la homología del
producto tensorial
\[
  M_{\bullet\bullet\bullet}
  =
  \bigotimes_{i=1}^{m}
  M_{\bullet\bullet\bullet}^{(i)}.
\]
Consideremos $M_{\bullet\bullet\bullet}$ como un complejo simple
$N_\bullet$, producto tensorial de los dos complejos simples
\[
  N_\bullet'
  =
  \bigotimes_{i=1}^{r}
  M_{\bullet\bullet\bullet}^{(i)},
  \qquad
  N_\bullet''
  =
  \bigotimes_{i=r+1}^{m}
  M_{\bullet\bullet\bullet}^{(i)},
\]
donde $N_\bullet'$ y $N_\bullet''$ están graduados por las sumas
de los grados totales de los $M_{\bullet\bullet\bullet}^{(i)}$.

Los $A$-módulos que forman $N_\bullet'$ y $N_\bullet''$ son
proyectivos. Por tanto, \sref{III.6.5.9} da
\[
  H_\bullet(M_{\bullet\bullet\bullet})
  =
  \mathbf{Tor}_\bullet^A(N_\bullet',N_\bullet'').
\]
La sucesión espectral buscada es, pues, precisamente
\sref{III.6.3.2.2} aplicada a $N_\bullet'$ y $N_\bullet''$, con
la interpretación precedente de sus módulos de homología para los
dos productos parciales
\[
  X^{(1)}\times_S\cdots\times_SX^{(r)}
  \quad\hbox{y}\quad
  X^{(r+1)}\times_S\cdots\times_SX^{(m)}.
\]
Finalmente, las afirmaciones de regularidad se deducen de
\sref{III.6.3.2}. En efecto, como los
$\sh P_\bullet^{(i)}$ están acotados inferiormente, también lo
están los $M_{\bullet\bullet\bullet}^{(i)}$ y, por consiguiente,
$N_\bullet'$ y $N_\bullet''$.
\end{proof}

\oldpage[III]{166}

